% Generated draft: requires mathematical and source review.
\documentclass[11pt]{article}
\usepackage[margin=25mm]{geometry}
\usepackage{fontspec}
\setmainfont{Latin Modern Roman}
\usepackage{amsmath,amssymb,mathtools,unicode-math}
\setmathfont{Latin Modern Math}
\usepackage{enumitem,needspace,xurl,hyperref}
\setlength{\parindent}{0pt}
\setlength{\parskip}{5pt}
\setlength{\emergencystretch}{3em}
\newcommand{\N}{\mathbb{N}}
\newcommand{\Z}{\mathbb{Z}}
\newcommand{\Q}{\mathbb{Q}}
\newcommand{\R}{\mathbb{R}}
\newcommand{\C}{\mathbb{C}}
\newcommand{\F}{\mathbb{F}}
\newcommand{\A}{\mathbb{A}}
\newcommand{\PP}{\mathbb{P}}
\newcommand{\E}{\mathbb{E}}
\newcommand{\eps}{\varepsilon}
\newcommand{\dd}{\,\mathrm{d}}
\begin{document}
\section*{UnsolvedMath ID: 30003197}
\textbf{Source problem code:} catalogue entry
\subsection{Definitions and mathematical statement}
Let \(S\) be a smooth quasi-projective variety and let \(V\) be a variation of mixed Hodge structures over \(S\), as in the supplied statement. Write \(\operatorname{HL}(S,V)\subseteq S\) for the Hodge locus intended by the source, subject to the unresolved conventions below. Define its Zariski closure by
\[
Z(S,V):=\overline{\operatorname{HL}(S,V)}^{\mathrm{Zar}}
=\bigcap_{\substack{Y\subseteq S\text{ Zariski closed}\\ \operatorname{HL}(S,V)\subseteq Y}}Y.
\]
Thus the Hodge locus is Zariski dense in \(S\) precisely when \(Z(S,V)=S\).

The requested problem is to determine, for the pairs \((S,V)\) under consideration, the geometric restrictions on \(Z(S,V)\), and to characterize those pairs for which
\[
\overline{\operatorname{HL}(S,V)}^{\mathrm{Zar}}=S.
\]
The source asks for a characterization and does not propose a particular necessary or sufficient condition.

The meaning of the Hodge locus requires further specification. Under a rational mixed Hodge structure convention, a Hodge class in a fiber \(H\) means a morphism \(\Q(0)\to H\) of mixed Hodge structures, where \(\Q(0)\) is the one-dimensional rational Hodge structure of weight zero and type \((0,0)\). A Hodge locus can refer to loci where locally flat classes become Hodge classes, possibly in specified tensor constructions and Tate twists of \(V\). The record does not specify those constructions, the coefficient lattice or field, or whether classes that persist generically are excluded. These choices are material: including the zero class without an exceptional-locus convention would give all of \(S\).

The usual setting involves a complex algebraic variety and a variation on its associated analytic space, but the record does not explicitly state the base field. It also does not specify admissibility, polarizability of the weight-graded pieces, or connectedness or irreducibility of \(S\). These hypotheses and the precise definition of \(\operatorname{HL}(S,V)\) must be fixed before the requested characterization has a unique mathematical meaning; none is silently imposed here.
\subsection{Short English statement}
For a variation of mixed Hodge structures \(V\) on a smooth quasi-projective variety \(S\), describe the possible Zariski closures of its Hodge locus and characterize when that locus is Zariski dense in \(S\). The source leaves the Hodge-locus convention and several hypotheses unspecified, so this is a statement-only draft requiring mathematical review.
\subsection{Sources}
\begin{itemize}
\item Local UnsolvedMath registry: \texttt{data/unsolvedmath/problems.json}, with \texttt{catalog/problems.json} designated as the supplementary registry. Permanent UnsolvedMath ID: \(30003197\). Supplied source code: \texttt{catalogue entry}. Title: \emph{Density and Closure of Mixed Hodge Loci}. The supplied record attributes the problem to Oberwolfach Reports workshop discussions; no individual proposer is identified, and the proposed year is recorded as 2016.
\item Supplied report citation: \emph{Arithmetic Geometry (2017)}, Oberwolfach Reports, DOI \texttt{10.4171/owr/2016/38}, pp.\ 2171--2224. \url{https://doi.org/10.4171/owr/2016/38}. The citation is reproduced as supplied; the relation between its 2017 designation and the 2016 DOI and proposed year is unresolved.
\item Supplied literature reference: Khelifa, \emph{On the distribution of mixed Hodge loci}. \url{https://arxiv.org/abs/2603.20272}.
\item Supplied literature reference: Klingler, \emph{Hodge locus conjectures}. \url{https://javier.fresan.perso.math.cnrs.fr/klingler.pdf}.
\item Literature assessment dated exactly \texttt{2026-08-21}: the supplied curation labels the problem partially solved and reports known Hodge-locus algebraicity and substantial distribution and density results, while reporting no general characterization of closures or dense cases. This is supplied dated curation derived from the OWR record, not an independently verified literature assessment. Its time-sensitive status and cited claims require independent verification.
\item The supplied statement status is \texttt{corrected\_verified}; the accompanying explanation says that mathematical content was preserved while notation, typography, or omitted definitions were repaired. This is a supplied editorial designation, not mathematical verification of this draft or independent verification of the report.
\end{itemize}
\end{document}
